\worldchapter{Ausrüstung}{equipment}
\label{ch:ausruestung}

\begin{multicols*}{2}

Die Ausrüstung eines Charakters bestimmt, womit er kämpft, wie gut er Treffer abfängt und welche Hilfsmittel ihm in einer Szene zur Verfügung stehen.
Dieses Kapitel beschreibt die allgemeinen Regeln für Ausrüstung.
Die konkreten Waffen, Rüstungen, Schilde und Gegenstände stehen gesammelt in \cref{app:equipment}.

\section{Waffen und Angriffe}\label{sec:waffen}

Jede Waffe besitzt feste Spielwerte.
Diese Werte bestimmen, wie sie in einem Konflikt eingesetzt wird.

\begin{description}[style=nextline,leftmargin=0pt]
\item[\textbf{Schaden}]\index{Schaden} Der Grundschaden eines erfolgreichen Treffers.
\item[\textbf{Reichweite}] Die Distanzstufen, auf denen die Waffe eingesetzt werden kann.
\item[\textbf{Griff}] Ein Modifikator für den ersten Angriff einer Runde.
\item[\textbf{Eigenschaften}] Besondere Regeln der Waffe.
\end{description}

\subsection{Griff}\label{subsec:griff}

Der \textbf{Griff} gibt an, wie leicht eine Waffe im entscheidenden Moment geführt werden kann.
Der Modifikator gilt für den ersten Angriff in einer Runde mit dieser Waffe.

\begin{itemize}[leftmargin=*]
\item \textbf{G2}: +10 auf den ersten Angriff der Runde mit dieser Waffe
\item \textbf{G1}: kein Modifikator
\item \textbf{G0}: -10 auf den ersten Angriff der Runde mit dieser Waffe
\end{itemize}

\subsection{Nachladen und Werfen}

Einige Waffen müssen nach einem Schuss oder Wurf wieder einsatzbereit gemacht werden.
Dies geschieht nach den Eigenschaften der Waffe.

\begin{itemize}[leftmargin=*]
\item Waffen mit \textbf{Nachladen 1} benötigen 1 Aktion zum Nachladen.
\item Waffen mit \textbf{Nachladen 2} benötigen 2 Aktionen zum Nachladen. Diese können aufgeteilt oder am Stück eingesetzt werden.
\item Waffen mit \textbf{Wurf} können auf nah geworfen werden und müssen danach wiederbeschafft werden.
\end{itemize}

\section{Rüstung, Schilde und Deckung}\label{sec:rustung}

\textbf{Rüstung}\index{Rustung} besitzt drei Kernwerte.

\begin{description}[style=nextline,leftmargin=0pt]
\item[\textbf{Schutz}]\index{Schutz} Schutz reduziert den Schaden eines Treffers. Der finale Schaden ergibt sich aus Schaden minus Schutz, mindestens jedoch 0.
\item[\textbf{Last}] Last erschwert Bewegung und Heimlichkeit.
\item[\textbf{Versiegelung}]\index{Versiegelung} Versiegelung schützt vor einzelnen Zuständen wie Erschütterung oder Blutverlust.
\end{description}

Für \textbf{Last} gelten folgende Stufen:

\begin{itemize}[leftmargin=*]
\item \textbf{L0}: keine Bewegungslast
\item \textbf{L1}: -10 auf Schleichbewegung
\item \textbf{L2}: -10 auf Schleichbewegung und Sprint nur 1 Band
\end{itemize}

Für \textbf{Versiegelung} gelten folgende Stufen:

\begin{itemize}[leftmargin=*]
\item \textbf{S0}: keine Zusatzregel
\item \textbf{S1}: ein erstes \textbf{Erschüttert} oder \textbf{Blutend} pro Konflikt ignorieren
\item \textbf{S2}: wie S1 und zusätzlich einmal pro Szene ein weiteres \textbf{Erschüttert} ignorieren
\end{itemize}

Es gilt immer nur \textbf{ein} getragenes Rüstungsprofil, auch wenn mehrere Rüstungsteile getragen werden.

\subsection{Schilde}\label{sec:schilde}

\textbf{Schilde}\index{Schild} werden als Reaktion eingesetzt.
Eine Schildreaktion kostet immer 1 Aktion und gilt nur gegen den auslösenden Angriff.

\begin{itemize}[leftmargin=*]
\item \textbf{Faustschild}: +10 gegen einen Nahkampfangriff
\item \textbf{Rundschild}: +10 gegen einen Angriff
\item \textbf{Dreiecksschild}: +20 gegen einen Fernkampfangriff
\end{itemize}

\subsection{Deckung}\label{sec:deckung}

\textbf{Deckung}\index{Deckung} wirkt nur gegen Fernangriffe.
Es zählt immer nur die höchste vorhandene Deckung.

\begin{itemize}[leftmargin=*]
\item \textbf{Leichte Deckung}: +10 Verteidigung
\item \textbf{Starke Deckung}: +20 Verteidigung
\end{itemize}

\section{Ausrüstungseigenschaften}\label{sec:eigenschaften}

Ausrüstungseigenschaften verändern den normalen Ablauf einer Handlung oder ergänzen ihn um feste Zusatzeffekte.

\begin{description}[style=nextline,leftmargin=0pt]
\item[\textbf{Abfangen}] Mit gehaltener Aktion kann auf einen Gegner reagiert werden, der von nah in den Nahkampf eintritt.
\item[\textbf{Aufsetzen}] Wurde in derselben Runde gezielt oder aufgesetzt und das Ziel tritt in den Nahkampf ein, verursacht der Angriff +1 Schaden.
\item[\textbf{Ausgewogen}] Keine zusätzliche Regel.
\item[\textbf{Brutal}] Bei einem sauberen oder starken Treffer muss das Ziel Körper + Standhaftigkeit testen, sonst wird es \textbf{Erschüttert}.
\item[\textbf{Durchdringend}] Ignoriert 1 Punkt Schutz.
\item[\textbf{Feldtauglich}] Für Marsch, Wache und Reise gemacht. Die Spielleiterin kann in passenden Situationen kleine Nachteile durch Witterung, Schmutz oder langes Tragen entfallen lassen.
\item[\textbf{Kriegsgerät}] Auffällige militärische Ausrüstung. Sie ist in zivilen oder engen Umgebungen schwer zu verbergen und kann sofort Aufmerksamkeit erregen.
\item[\textbf{Laut}] Der Einsatz ist auffällig und kann sofort Aufmerksamkeit oder Eskalation auslösen.
\item[\textbf{Nichttödlich}] Verursachte Wunden können stattdessen als gleich hohe Bürde vergeben werden.
\item[\textbf{Ritual}] Das Stück ist für zeremonielle, gelehrte oder kultische Handlungen ausgelegt. Die Spielleiterin kann in passenden Situationen kleine Vorteile auf Vorbereitung, Anerkennung oder Deutung gewähren.
\item[\textbf{Schwer}] -10 auf Reaktionen mit diesem Ausrüstungsstück.
\item[\textbf{Verbergbar}] +10, um den Gegenstand vor Szenenbeginn zu verbergen.
\item[\textbf{Verstärkt}] Die Ausrüstung ist besonders belastbar. Die Spielleiterin kann in passenden Situationen Abnutzung, Bruch oder Beschädigung später eintreten lassen.
\item[\textbf{Wuchtig}] Gegen Schutz 2+ verursacht der Angriff +1 Schaden vor Abzug des Schutzes.
\item[\textbf{Nah}] +10 in engem Raum.
\item[\textbf{Nachladen 1}] Nachladen kostet 1 Aktion.
\item[\textbf{Nachladen 2}] Nachladen kostet 2 Aktionen, aufgeteilt oder am Stück.
\item[\textbf{Wurf}] Kann auf nah geworfen werden und muss danach wiederbeschafft werden.
\item[\textbf{Zweihändig}] Kann nicht zusammen mit einem Schild geführt werden.
\end{description}

\section{Gebrauchsgegenstände}\label{sec:gebrauchsgegenstande}

Gebrauchsgegenstände decken Werkzeuge, Heilmittel, Reisebedarf und ähnliche Dinge ab.
Wenn die Regeln nichts anderes angeben, kostet der Einsatz eines Gegenstands 1 Aktion.
Verbrauchsgüter sind in der Regel einmalig oder besitzen nur wenige Anwendungen.

Der vollständige Satz aus Waffen, Rüstungen, Schilden und Gebrauchsgegenständen steht in \cref{app:equipment}.
\end{multicols*}
